% 核心结论：对称性和边界条件共同决定适合的坐标系。
% 图契约：物理竞赛讲义；印刷宽16cm，最小字号8.5pt；几何示意与公式分区。
% 布局：[几何/曲线 7cm] [定义/结果 7cm]；实线主量，虚线辅助；无装饰数据。
\documentclass[tikz,border=6pt]{standalone}
\usepackage{fontspec,xeCJK,amsmath,amssymb,graphicx}
\IfFontExistsTF{Times New Roman}{\setmainfont{Times New Roman}}{\setmainfont{TeX Gyre Termes}}
\IfFontExistsTF{Heiti SC}{\setCJKmainfont[AutoFakeBold=1.5]{Heiti SC}}{\setCJKmainfont[AutoFakeBold=1.5]{FandolHei-Regular}}
\usetikzlibrary{arrows.meta,bending,calc,positioning,patterns}
\definecolor{ink}{HTML}{243745}
\definecolor{blue}{HTML}{275DAB}
\definecolor{green}{HTML}{177C63}
\definecolor{red}{HTML}{BD4935}
\definecolor{orange}{HTML}{AA6C16}
\definecolor{purple}{HTML}{715393}
\tikzset{
 every node/.style={font=\fontsize{9.5}{13}\selectfont,text=ink},
 axis/.style={-{Stealth[length=4pt]},line width=.65pt,draw=ink!65},
 vec/.style={-{Stealth[length=5pt]},line width=1pt,draw=blue},
 short/.style={-{Stealth[length=3pt]},line width=.8pt},
 guide/.style={dashed,draw=ink!40,line width=.5pt},
 card/.style={rounded corners=2pt,draw=blue!60,fill=blue!4,
   align=center,text width=6.8cm,inner xsep=9pt,inner ysep=7pt},
 title/.style={font=\fontsize{11}{14}\selectfont\bfseries,text width=15cm,align=center},
 label/.style={font=\fontsize{8.5}{11}\selectfont,align=center}}
\begin{document}
\begin{tikzpicture}[x=1cm,y=1cm]
\path[use as bounding box] (-8,-3.35) rectangle (8,3.35);
\node[title] at (0,2.95) {坐标选择：先看源与边界的对称性};

\node[card,text width=6cm] at (-4,1.2) {平面、矩形边界\\优先考虑直角坐标 $(x,y,z)$};
\node[card,text width=6cm,draw=green!65,fill=green!4] at (4,1.2) {无限长线源、同轴柱面\\优先考虑柱坐标 $(\rho,\phi,z)$};
\node[card,text width=6cm,draw=orange!65,fill=orange!4] at (-4,-.7) {点源、同心球形边界\\优先考虑球坐标 $(r,\theta,\phi)$};
\node[card,text width=6cm,draw=purple!65,fill=purple!4] at (4,-.7) {复杂曲面边界\\选能简化边界条件的正交坐标};
\node[label,text width=14cm] at (0,-2.35) {选定坐标后：检查对称性与定义域，再使用尺度因子、微元和算符公式。};

\end{tikzpicture}\end{document}
